\originalchapter{网络、集体运动与DNA包装}
\originalsection{网络结构和阈值}
用节点表示基因、蛋白或细胞，边表示已经定义的相互作用。度数、路径长度和聚类系数描述的是不同问题；仅凭度数分布不能推断一条边的因果作用。随机图中平均度 $c$ 的稀疏独立连边模型，探索邻居像平均子代数为 $c$ 的分枝过程。令 $u$ 为沿一条边进入的分枝最终有限的概率，则 $u=\e^{-c(1-u)}$。巨连通分量比例 $S=1-u$ 在 $c>1$ 才出现非零解。

更一般的无关联度分布 $P(k)$，沿边看到的度分布为 $kP(k)/\avg{k}$，可继续探索的平均新边数是 $\avg{k(k-1)}/\avg{k}$。因此阈值涉及二阶矩，而不只是平均度。这个论证假设局部近似树状；模块化、度相关和强聚类网络需要修正。

\originalsection{反应网络的稳定性}
浓度 $\mathbf c$ 的动力学 $\dot{\mathbf c}=\mathbf f(\mathbf c)$ 在稳态附近写成 $\dot{\delta\mathbf c}=J\delta\mathbf c$，其中 $J_{ij}=\partial f_i/\partial c_j$。所有特征值实部为负才保证局部渐近稳定。网络中的一条负反馈边并不单独决定稳定性，反馈强度、延迟和其他回路同样影响特征值。

以具有饱和合成和线性降解的调控模型 $\dot c=\alpha/[1+(c/K)^n]-\delta c$ 为例，右端导数为负，稳态唯一且一维扰动衰减；它本身不能持续振荡。若加入转录与翻译两级、时滞或额外状态，才可能产生振荡机制。课程网络主题的部分页面只有纲要，本节给的是编者独立模型说明，不代表取得了不存在的详细讲义。

\originalsection{包装能量与尺度选择}
噬菌体衣壳内的DNA密度高，热运动、弯曲和排斥共同决定包装功。若长度 $L$ 主要弯成半径 $R$ 的环，弯曲能尺度 $E_b\sim\kb T\ell_pL/(2R^2)$；若曲率随位置变化，应实际积分 $\kappa\int R(s)^{-2}\dd s/2$，不能以最外层半径代表全部内圈。线电荷的屏蔽排斥还需要盐浓度、有效电荷和相邻链距离。

包装马达的力是相应自由能对已装入长度的导数，而不是总能量除以任意长度。随着内部空隙缩小，弯曲半径下降，力常迅速增大；由平衡自由能算出的准静态力不包含黏性耗散。将一种模型拟合得好也不表明所有能量项已唯一确定。关于衣壳包装的原题需要额外材料几何和实验资料，本版保留课程来源，不把它列为已完成题解。
